\chapter{What ProPrep is (and is not)}
\label{ch:whatis}

ProPrep helps you turn a protein structure from a database into a system that
is ready to simulate with the Amber programs. It walks you through every step,
asks you a short question at each one, and explains why the question matters.
This chapter explains what that means in plain words, what ProPrep does for you,
and what it leaves to you.

\section{The problem ProPrep solves}
\label{sec:whatis-problem}

A structure from the \term{Protein Data Bank} (PDB, the public archive of
experimentally solved structures) is a picture, not a model ready for
simulation. A typical file has:
\begin{itemize}
  \item no hydrogen atoms, because X-ray experiments rarely see them;
  \item missing pieces, such as loops or side chains that were too mobile to
        see;
  \item extra molecules: crystal waters, salts, and additives from the
        experiment;
  \item groups the simulation programs do not recognize on their own, such as
        a heme, a metal ion, or a drug molecule;
  \item amino acids whose charge depends on the \term{pH} (the acidity of the
        surroundings), such as histidine.
\end{itemize}

Before a simulation can start, someone must decide what to keep, fix what is
missing, choose the charge of every titratable residue, give every unusual
group its parameters, and put the protein in a box of water and salt. Doing this
by hand means many small programs, file formats and naming rules. ProPrep does
that labor for you and keeps a record of every choice.

\section{Words you will meet}
\label{sec:whatis-words}

\begin{description}
  \item[A prepared structure] is a complete, consistently named model of your
        protein. It has the parts you decided to keep, the missing atoms filled
        in, and residue names that match the charge states you chose.
  \item[A molecular dynamics simulation] (\term{MD}) is a calculation that
        moves every atom a tiny step at a time, using the forces between the
        atoms. Millions of steps make a short movie of how the protein moves in
        water. The movie is saved as a \term{trajectory} file.
  \item[A force field] is the set of numbers that describes these forces:
        bond lengths, angles, charges, and so on. ProPrep uses the Amber
        \term{force field}s, for example ff19SB for proteins.
  \item[A topology] is the file that lists every atom of your system with its
        force-field numbers. Together with a \term{coordinate file} (where each
        atom is), it is what the simulation program reads.
\end{description}

\section{What ProPrep does}
\label{sec:whatis-does}

ProPrep covers the whole path from a PDB code to an analyzed simulation. The
main stages, each with its own chapter in Part~IV, are:

\begin{enumerate}
  \item \textbf{Get a structure} from the PDB, from the AlphaFold database, or
        from your own file (Chapter~\ref{ch:structure}).
  \item \textbf{Clean and repair it}: choose which chains, waters and ligands
        to keep, fill gaps, and make mutations (Chapter~\ref{ch:cleaning}).
  \item \textbf{Set protonation states} at your chosen pH with PROPKA, with a
        reason shown for each residue (Chapter~\ref{ch:protonation}).
  \item \textbf{Handle special sites}: hemes, iron--sulfur clusters, metals
        and ligands (Chapter~\ref{ch:redox} and Chapter~\ref{ch:forcefields}).
  \item \textbf{Build the system}: force field, water box and salt, written as
        an editable input file for Amber's tLEaP program
        (Chapter~\ref{ch:topology}).
  \item \textbf{Run and analyze the simulation} on your computer or on a
        computing cluster (Chapters~\ref{ch:runningmd} and~\ref{ch:analysis}).
\end{enumerate}

It also builds membranes around membrane proteins (Chapter~\ref{ch:membrane}),
docks small molecules (Chapter~\ref{ch:docking}), writes input files for
QM/MM calculations (Chapter~\ref{ch:qmmm}), and compares and shows structures
in 3D (Chapter~\ref{ch:compare}).

\section{It automates the labor, not the decisions}
\label{sec:whatis-idea}

The idea behind ProPrep is stated in its own description: it \emph{automates
the labor of preparation, not the decisions}. Each step is a small choice,
usually a single key, shown together with what the step does, why it matters
and how ProPrep will carry it out. You stay informed and in control instead of
accepting hidden settings.

In practice this means:
\begin{itemize}
  \item ProPrep \textbf{suggests} a choice (the default in brackets), but you
        make it. Pressing \key{Enter} accepts the suggestion.
  \item ProPrep \textbf{explains} each scientific choice on the screen, often in
        a box with a short background text.
  \item ProPrep \textbf{records} every answer in a \term{session file}. You
        can replay the file, correct a wrong answer, or use it as a recipe for
        many proteins (Chapters~\ref{ch:sessions} and~\ref{ch:batch}).
  \item ProPrep \textbf{shows} the files it writes, such as the tLEaP input, so
        you can read and edit them before they are used.
\end{itemize}

ProPrep works with text: menus, questions and answers. This keeps it usable
over a remote connection to a computing cluster, and for researchers who use a
screen reader. A 3D structure viewer in the web browser is available, but you
never need it.

\section{What ProPrep does not do}
\label{sec:whatis-not}

\begin{itemize}
  \item \textbf{It does not make your scientific decisions.} It will not tell
        you which pH, which chains or which water model is right for your
        question. In our tests it also accepted some combinations that its
        own lists advise against (for example the TIP3P water model with
        ff19SB) without a warning. Read the explanations and decide.
  \item \textbf{It does not run quantum-chemistry programs for you.} To make
        new parameters, and for QM/MM, it writes input files for Gaussian or
        ORCA. You need your own copy of those programs to run them
        (Appendix~\ref{app:software}).
  \item \textbf{It is not a simulation engine.} It drives Amber's programs.
        The free installation includes \texttt{sander}; the faster
        \texttt{pmemd} is a separate Amber program that you install
        yourself.
  \item \textbf{It prepares systems for Amber only.} The files it writes are
        Amber files (\file{.prmtop}, \file{.rst7}, \file{.mdin}).
  \item \textbf{It does not replace checking.} ProPrep reports what it did, but
        a wrong choice upstream gives a wrong system downstream. Look at the
        structure and the log files before you trust a long simulation.
\end{itemize}

\section{Who makes ProPrep, and how to cite it}
\label{sec:whatis-cite}

ProPrep is developed by the Guberman-Pfeffer Lab at New Mexico State
University, with contributions from Adara Walker and John Cairns. It is free
and open source under the MIT license. The source code is at
\url{https://github.com/mjgplab/proprep}.

When you leave ProPrep, it shows how to cite it:

\begin{screen}
┌───────────────────── Thank you for using ProPrep ──────────────────────┐
│ If you found ProPrep useful, please cite:                              │
│                                                                        │
│ Walker, A. & Guberman-Pfeffer, M.J. (2026). ProPrep: An Interactive    │
│ and Instructional Interface for Proper Protein Preparation with AMBER. │
│ bioRxiv. https://doi.org/10.64898/2026.02.26.708365                    │
│                                                                        │
│ Feedback (positive or negative) is welcome at: mjgp@nmsu.edu           │
└────────────────────────────────────────────────────────────────────────┘
\end{screen}

The developers ask you to cite both the article and the software itself. Each
software release is archived with a \term{DOI} (a permanent web address for a
publication or data set): \url{https://doi.org/10.5281/zenodo.21764540}. The
file \file{CITATION.cff} in the ProPrep source code has the full details.

\begin{tip}
Write down the ProPrep version you used (\texttt{proprep --version} prints it)
in your lab notebook and in your methods section. Session files are tied to the
version that recorded them.
\end{tip}
